\documentclass{report}
\usepackage{amsmath,amssymb}
\setlength{\parindent}{0mm}
\setlength{\parskip}{1em}
\begin{document}
\begin{center}
\rule{6in}{1pt} \
{\large Stephanie Friedhoff \\
{\bf A multigrid-in-time algorithm for solving evolution equations in parallel}}

Tufts University \\ Department of Mathematics \\ 503 Boston Avenue \\ Medford MA 02155
\\
{\tt Stephanie.Friedhoff@tufts.edu}\\
Robert Falgout\\
Tzanio Kolev\\
Scott MacLachlan\\
Jacob Schroder\end{center}

We consider optimal-scaling multigrid solvers for the linear systems that
arise from the discretization of problems with evolutionary behavior.
Typically, solution algorithms for evolution equations are based on a
time-marching approach, meaning solving for one time step after the
other. These traditional time integration techniques lead to
optimal-scaling but not parallelizable algorithms. However, current
trends in computer architectures are leading towards more, but slower,
processors and, therefore, driving the need for greater parallelism. One
approach to achieve parallelism in time is with multigrid, but while
classical multigrid methods rely on multiscale representations in space,
that arise naturally from decomposing a function into a hierarchy of
frequencies from global smooth modes to local oscillations, these
approaches do not extend to evolutionary variables in a straightforward
manner, because of the fundamentally local structure of the evolution. In
this talk, we present an optimal and scalable multigrid-in-time algorithm
for diffusion equations as simple examples of evolution equations. Our
algorithm is based on interpreting the parareal time integration method
[J. L. Lions, Y. Maday, and G. Turinici, R\'{e}solution d'EDP par un
sch\'{e}ma en temps parar\'{e}el, C.R. Acad. Sci. Paris, S\'{e}r. I Math,
332 (2001), pp. 661-668] as a two-level reduction scheme, and developing
a multilevel algorithm from this viewpoint. We demonstrate optimality of
our algorithm for solving the one-dimensional diffusion equation in
numerical experiments. Furthermore, by using parallel performance models,
we show that we can expect speedup in comparison to sequential
time-marching on modern architectures.


\end{document}
